\documentclass{article}
\usepackage{amsmath}
\usepackage[repeqn=same]{../../apxproof}
\usepackage{hyperref}

\newtheoremrep{theorem}{Theorem}
\newtheoremrep{lemma}[theorem]{Lemma}
\newtheoremrep*{claim}[theorem]{Claim}

\begin{document}

\section{Main text}

% Repeated theorems nested in one another: the inner ones are restated in
% the appendix as part of the outer one, not on their own.
\begin{theoremrep}[Outer]\label{th:outer}%
  An outer theorem,
  \begin{lemmarep}[Inner]\label{lem:inner}%
    with an inner lemma
    \begin{equation}\label{eq:inner}1=1\end{equation}
    \begin{theoremrep}\label{th:innermost}%
      containing a theorem of the same kind,
    \end{theoremrep}
  \end{lemmarep}
  a starred claim
  \begin{claimrep}\label{cl:inner}%
    that is restated with it,
  \end{claimrep}
  and some text after them.
\end{theoremrep}

\begin{proof}
  Proof of the outer theorem.
\end{proof}

\begin{theoremrep}\label{th:after}%
  A theorem after the nested ones.
  \begin{equation}\label{eq:after}2=2\end{equation}
\end{theoremrep}

\begin{proof}
  Proof of the theorem after.
\end{proof}

References: Theorem~\ref{th:outer}, Lemma~\ref{lem:inner},
Theorem~\ref{th:innermost}, Claim~\ref{cl:inner}, Theorem~\ref{th:after},
equations~\eqref{eq:inner} and~\eqref{eq:after}, and the \texttt{-apx}
variants \ref{th:outer-apx}, \ref{lem:inner-apx}, \ref{th:innermost-apx},
\ref{cl:inner-apx}, \ref{th:after-apx}, \eqref{eq:inner-apx},
\eqref{eq:after-apx}.

\end{document}
